\begin{python0}
from solutions import *; clear()
\end{python0}
\sectionthree{Static members variables (in a class)}

Besides static variables in a code block, classes can have static
members as well. Let me give the concepts first.

(The code in this section is not runnable. You'll have
to wait for the next section before I complete the code.)

Think of \textbf{static member variables} as \textbf{variables belonging
to a class} and not to each object. There is only one copy for each
static member. All objects of a class \textbf{share} the same static
member. If \texttt{y} is a static member of a class and \texttt{obj1} and
\texttt{obj2} are objects of that class, then \texttt{obj1.y} is the
\textbf{same} as \texttt{obj2.y}.

To make a member static put the keyword \texttt{static} in front of the
declaration of the member in the class declaration. Static members can
be public or private.

For the next example, I'm going to make all members
public just to simplify the discussion. Don't run the
code yet -- I'm going through the concept of static
members and the code won't be runnable until the next
section when I show you how to initialize static members.

First let me show you how to create static member variables.
Here's class C:

\begin{consolethree}[escapeinside=||]
// C.h
#ifndef C_H
#define C_H

class C
{
public:
    int x;               // regular class member variable
    |{\bf static}| int y;       // static class member variable
    |{\bf static}| int z;       // static class member variable
};

#endif
\end{consolethree}

Look at the ``\texttt{static}'' words in the code. This class has two
static members: \texttt{y} and \texttt{z}.

Now if I do this in \texttt{main()} (remember: don't run
the code yet -- just focus on the idea):

\begin{consolethree}[escapeinside=||]
#include <iostream>
#include "C.h"

int main()
{
    C c, d, e; // 3 objects from class C
    ...
\end{consolethree}

then the objects currently in memory and the static members of class
\texttt{C} look like this (conceptually):

Note the following:

\begin{itemize}
\item Each of the 3 objects have their own member variable x. So ...
  \texttt{c.x} is different from \texttt{d.x} which is different from
  \texttt{e.x}. When c changes its x to \texttt{1000}, \texttt{d.x} is still
  \texttt{1} and \texttt{e.x} is still \texttt{2}.
\item All 3 objects \textbf{share} \texttt{y} and \texttt{z} (the
  \textbf{static} member variables). All 3 objects do have access to
  \texttt{y} and \texttt{z}. So when \texttt{c} changes \texttt{y} to \texttt{43}
  and \texttt{d} print \texttt{y}, \texttt{d} will be printing \texttt{43}.
\end{itemize}

Get it? It's really not that complex: regular member
variables are not shared among objects while static member variables are
shared across all objects.

Now although the objects \texttt{c}, \texttt{d}, \texttt{e} have access to the
static members \texttt{y}, \texttt{z}, you should view the static member
variables as variables belonging to the class \texttt{C}. So ...
although you can write \texttt{c.y} in your code (member \texttt{y} is
public -- look at definition of class \texttt{C}):

\begin{consolethree}[escapeinside=||]
#include <iostream>
#include "C.h"

int main()
{
    C c, d, e; // 3 objects from class C
    std::cout << c.y << '\n';
    ...
\end{consolethree}

it's more common to write \texttt{C::y}.

\begin{consolethree}[escapeinside=||]
#include <iostream>
#include "C.h"

int main()
{
    C c, d, e; // 3 objects from class C
    std::cout << |{\bf C::y}| << '\n'; // better than c.y
    ...
\end{consolethree}

If y is a private member:

\begin{consolethree}[escapeinside=||]
...
class C
{
public:
    int x;               // regular class member variable
    static int z;        // static class member variable
private:
    static int y;        // static class member variable
};
...
\end{consolethree}

then this won't work:

\begin{consolethree}[escapeinside=||]
#include <iostream>
#include "C.h"

int main()
{
    C c, d, e; // 3 objects from class C
    std::cout << C::y << '\n'; // BAD! BAD! BAD!
    ...
\end{consolethree}

Now let me show you how to initialize the value of static member
variables so that you can actually run the above code.


